%! Author = markt
%! Date = 3/31/2025

% Preamble
\documentclass[11pt]{article}

% Packages
\usepackage{latexsym,amsfonts,amssymb,amsthm,amsmath, yfonts, cancel, listings, xcolor, graphicx, article}
\usepackage{color}
\usepackage{listings}

\author{Mark Stanley}

\definecolor{codegreen}{rgb}{0,0.6,0}
\definecolor{codegray}{rgb}{0.5,0.5,0.5}
\definecolor{codepurple}{rgb}{0.58,0,0.82}
\definecolor{backcolour}{rgb}{0.95,0.95,0.92}
\lstdefinestyle{mystyle}{
    backgroundcolor=\color{backcolour},
    commentstyle=\color{codegreen},
    keywordstyle=\color{magenta},
    numberstyle=\tiny\color{codegray},
    stringstyle=\color{codepurple},
    basicstyle=\ttfamily\footnotesize,
    breakatwhitespace=false,
    breaklines=true,
    captionpos=b,
    keepspaces=true,
    numbers=left,
    numbersep=5pt,
    showspaces=false,
    showstringspaces=false,
    showtabs=false,
    tabsize=2
}

\title{Math 443 Review Session}

% Document
\begin{document}
    \lstset{style=mystyle}

    \maketitle

    \section{Gaussian Elimination}
    Create an upper triangular matrix from a given matrix.
    \[A=\begin{pmatrix}1 & 2 \\ 4 & 5\end{pmatrix}\]
    In Gaussian Elimination, you are allowed to add or subtract a scalar value of any row to another row.
    So, for example here, we would subtract 4 times the first row from the second row. $R_2 \leftarrow R_2 - 4R_1$

    \section{Elementary Matricies}
    An elementary matrix corresponds to a single row operation applied to the identity matrix. Such as
    1.
    Adding a scalar multiple of one row to another row.
    This is just the identity matrix with the scalar multiple in the row that is being added to the other row.
    2.
    Swapping two rows.
    This is just the identity matrix with the two rows swapped.
    3.
    Multiplying a row by a nonzero scalar.
    This is just the identity matrix with the row multiplied by the scalar.

    \textbf{Example:}
    For the operation we applied earlier, we can represent it as a matrix.
    \[E = \begin{pmatrix}1 & 0 \\ -4 & 1\end{pmatrix}\]
    Then, we can multiply this matrix by the original matrix to get the upper triangular matrix.

    \section{Inverse Elementary Matricies}
    To reverse the operation of an elementary matrix, we can just multiply by the inverse of the elementary matrix.
    For the scalar addition of rows, the inverse is simply the negative of the scalar from the $(j, i)$ position put in
    the $(i, j)$ position.

    \textbf{LU Factorization:}
    Given a matrix $A$, we can factor it into the product of a lower triangular matrix $L$ and an upper triangular matrix $U$.
    Where $LU = (L_1L_2\ldotsL_n)(E_n\ldots E_2E_1A)=A$.
    So, simply perform Gaussian Elimination on $A$ and keep track of the elementary matrices.
    This will give you $L$ and $U$.

    \tectbf{Theorem:}
    A matrix $A$ is regular iff $A$ can be factored into the product of a lower triangular matrix $L$ and an upper triangular matrix $U$.

    \section{Gauss-Jordan Elimination:}
    The same principle as Gaussian Elimination, but we also eliminate the entries above the diagonal.
    This will give us the identity matrix.
    Used to solve systems of equations ($Ax=b$)

    \section{LDV Decomposition:}
    Given a matrix $A$, we can factor it into the product of a lower triangular matrix $L$, a diagonal matrix $D$, and an upper triangular matrix $V$.
    Where $LDV = A$.
    Similar to LU Decomposition.
    Perform Gaussian Elimination on $A$ and keep track of the elementary matrices.
    This will give you $L$ and $U$.
    $D$ will just be the diagonal of $U$.
    Then, divide each row of $U$ by the corresponding entry in $D$ to get $V$.

    \section{Determinant:}
    det$(A)$ is a function that returns a scalar value from a given matrix $A$.
    If det$(A) \neq 0$, then $A$ is invertible (regular).
    Else if det$(A) = 0$, then $A$ is singular (irregular).

    \textbf{Properties of Determinants:}
    1.
    Adding a scalar multiple of one row to another row does not change the determinant.
    2.
    Swapping two rows negates the determinant.
    3.
    Multiplying a row by a scalar multiplies the determinant by that scalar.
    4.
    The determinant of an upper triangular matrix is the product of the diagonal entries.

    To determine the determinant of a matrix, perform Gaussian Elimination to get the upper triangular matrix.
    Then, multiply the diagonal entries of the upper triangular matrix.
    If you swapped two rows, negate the determinant that many times.


    \textbf{Vector Spaces:}
    A vector space $V$ is a set of vectors that is closed under addition and scalar multiplication (i.e. $v,w \in V \quad v+w \in V$ and $v \in V \quad \alpha v \in V$).

    \textbf{Axioms:}
    1.
    $v+w = w+v$
    2.
    $(v+w)+u = v+(w+u)$
    3.
    There exists a zero vector $0$ such that $v+0=v$
    4.
    For every $v$, there exists a vector $-v$ such that $v+(-v)=0$
    5.
    $\alpha(v+w) = \alpha v + \alpha w$
    6.
    $\alpha(\beta v)=(\alpha \beta )v$
    7.
    $1v = v$

    \textbf{Span and Linear Independence:}
    The span of a set of vectors is the set of all possible linear combinations of those vectors.
    \[\textrm{Span}(v_1, \ldots, v_k)=\{c_1 v_1 + c_2 v_2 + \cdots + c_k v_k | c_i \in \mathbb{R}\}\}\]

    A set of vectors in linearly independent if the only linear combination of those vectors that equals the zero vector is the trivial one.
    In other words: \\
    Vectors $v_1, v_2, \ldots, v_k$ are linearly independent iff the following equation
    \[c_1 v_1 + c_2 v_2 + \cdots + c_k v_k = 0\]
    only has the trivial solution $c_1 = c_2 = \cdots = c_k = 0$.


    \textbf{Four Fundamental Subspaces:}
    1.
    The column space: The span of the columns of a matrix.
    2.
    The row space: The span of the rows of a matrix.
    3.
    The null space: The set of all solutions to $Ax=0$.
    4.
    The left null space: The set of all solutions to $A^Ty=0$.

    \textbf{Rank Nullity Theorem:}
    \[\textrm{rank}(A)+\textrm{nullity}(A)=n\]
    \[\textrm{rank}(A^T)+\textrm{nullity}(A^T)=m\]

    \section{Subspaces:}
    A subspace is a subset of a vector space that is itself a vector space.
    It must satisfy the axioms of a vector space, as well as have the same zero vector as the original vector space.

    \section(Inner Product:)
    \textbf{Definition:} An inner product of a vector space $V$ is a pairing $\langle v, w, \rangle \in mathbb{R}$ satisfying:
    1.
    \textbf{Bilinearity:} $\langle cu+dv, q \rangle = c\langle u, w \rangle + d \langle v , w \rangle $
    2.
    \textbf{Symmetry:} $\langle u, v \rangle = \langle v, u \rangle$
    3.
    Positivity: $\langle v, v \rangle \geq 0$ for $v \neq 0$, and $\langle 0, 0 \rangle = 0$

    \textbf{Examples of Inner Products:}
    1.
    The dot product: $\langle u, v \rangle = u \cdot v = \sum v_i w_i$
    2.
    The weighted dot product: $\langle u, v \rangle = \sum c_u v_i w_i$, where $c_i > 0$

    \section{Cauchy-Schwarz Inequality:}\label{sec:cauchy-schwarz}
    \textbf{Statement:} For all vectors $v, w$ in an inner product space,
    \[|\langle v, w \rangle| \leq ||v|| \cdot ||w||\]
    Equality holds iff $v$ and $w$ are linearly dependent.

    \section{Triangle Inequality and Minowski Inequality:}\label{sec:triangle}
    \textbf{Triangle Inequality:} For all vectors $v, w \in V$:
    \[||v+w|| \leq ||v|| + ||w||\]
    \textbf{Minowski Inequality:} For all vectors $v, w \in V$ and all $p \geq 1$:
    \[(\sum |v_i + w_i|^p)^{\frac{1}{p}} \leq (\sum |v_i|^p)^{\frac{1}{p}} + (\sum |w_i|^p)^{\frac{1}{p}}\]

    \section{Orthogonal Matrix:}
    A square matrix $Q$ is orthogonal if its rows and columns are orthogonal unit vectors.
    Meaning: for $c_1, \ldots, c_n \in Q$, $\|c_i\|=1$ and $c_i \cdot c_j = 0$ for $i \neq j$.
    or:\\
    \[Q^{TQ=QQ^T=I}\]

    \textbf{Key Properties of Orthogonal Matrices:}
    1.
    \textbf{Inverse equals transpose:} $Q^{-1}=Q^T$
    2.
    \textbf{Preserves length:} $||Qv||=||v||$
    3.
    \textbf{Determinant:} $|\det(Q)|= \pm 1$
    4.
    \textbf{Rotation and reflection:} Orthogonal matrices represent rotations and reflections.

    \section{Positive Definite Matrices:}
    A symmetric matrix $A$ is positive definite if for all nonzero vectors $v$, $v^TAv > 0$.

    \subsection*{(a)} Suppose $A$ is a positive definite matrix.
    Show that $A^{-1}$ is also positive definite.

    \begin{proof}
        We know that $x^T A x > 0, \, \forall x \in V$.
        Assume $y$ is any vector, we can define the product as $x = A^{-1} y$.
        We can now substitute this into the equation:
        \[y^T A^{-1} y > 0\]
        Thus, $A^{-1}$ is also positive definite.
    \end{proof}






    \begin{proof}
        Some Proof
    \end{proof}

    Some code:
    \begin{lstlisting}[language=Python]
    def python_code():
        pass
    \end{lstlisting}

\end{document}